\section{Finite based lists on \texorpdfstring{$p$}{p}-subgroups}\label{sec:sylow}

Exact comparison and fixed-total-order rigidity now meet.  We first
extend the comparison to a Morita equivalence of crossed orders.  We
then realize the relevant orders as actual blocks and apply the
numerical finiteness criterion.  Finally we recover the labelled
components and the identity markings using
Proposition~\ref{prop:marked}.

\subsection{Extending an exact comparison}

The following is the one-term crossed-order version of the
diagonal extension argument of Marcus
\cite[Theorem~3.4]{Marcus1996}, also recorded in
\cite[Lemma~10.2.8]{Rouquier1998}.  We give its algebraic
proof to retain the specified multiplication factors.

\begin{lemma}[Extension of a comparison bimodule]\label{lem:induced}
Let $B,B'$ be $\cO$-orders with normalized crossed systems
$(\alpha_s,a_{s,t})$ and $(\alpha'_s,a'_{s,t})$ for a finite group
$S$, and let $A,A'$ be their crossed orders.  Suppose an
$(B',B)$-Morita bimodule $\mathcal M$ has invertible
$\cO$-linear maps $J_s$ satisfying
Equations~\eqref{eq:comparison-covariance} and
\eqref{eq:comparison-exact}.  Then $A$ and $A'$ are
$\cO$-linearly Morita equivalent.
\end{lemma}

\begin{proof}
The induced right module carries the required left crossed action
precisely because the comparison has the exact factor identity.
Put $P=\mathcal M\otimes_B A$, a right $A$-module.  The left
$B'$-action is the original action on $\mathcal M$.  Define
right $A$-linear operators
\[
 L_s(v\otimes a)=J_s(v)\otimes u_s a.
\]
They are well defined on the balanced tensor product: covariance
and $u_sb=\alpha_s(b)u_s$ identify the images of
$vb\otimes a$ and $v\otimes ba$.  They are invertible, and
covariance gives $L_s b'=\alpha'_s(b')L_s$.
The exact factor identity gives
\begin{align*}
 L_sL_t(v\otimes a)
 &=a'_{s,t}J_{st}(v)a_{s,t}^{-1}\otimes
                 a_{s,t}u_{st}a\\
 &=a'_{s,t}L_{st}(v\otimes a).
\end{align*}
Thus these operators define a left $A'$-action.

The right $B$-module $\mathcal M$ is a finitely generated
projective generator.  Inducing its projective summand and
generator identities to $A$ shows that $P$ is a projective
generator as a right $A$-module.  It remains to identify its
endomorphism order with $A'$.

As a right $B$-module, $P$ is the direct sum of the components
$\mathcal M\otimes_B Bu_s$.  Induction--restriction adjunction
gives the following isomorphisms of $\cO$-modules:
\[
 \End_A(P)\cong\Hom_B(\mathcal M,P)
 \cong\bigoplus_{s\in S}
 \Hom_B(\mathcal M,\mathcal M\otimes_B Bu_s).
\]
Explicitly, a $B$-linear map $f$ extends uniquely to the
$A$-endomorphism $v\otimes a\mapsto f(v)a$.  For a map in
the degree-$s$ summand, composing its extension with $L_s^{-1}$
gives a degree-one endomorphism, hence an element of
$\End_B(\mathcal M)=B'$.  Therefore
\[
 \End_A(P)=\bigoplus_{s\in S}B'L_s.
\]
The relations already checked identify this order with $A'$.
The projective-generator characterization of Morita equivalence
now proves the lemma.
\end{proof}

\subsection{The actual extension blocks}

\begin{proposition}[Sylow restrictions are block orders]
\label{prop:sylow-block}
For the data of Proposition~\ref{prop:structure}, let
$S\leq Q\cap X$ be a pure $p$-subgroup.  The restriction of the
crossed order to $S$, with identity component $B_0$, is a block
of a finite group.  Its defect order is bounded in terms of
$p,M$, and its integral Morita--Frobenius number is bounded in
the same parameters.  All these restriction blocks have
finitely many integral Morita classes.
\end{proposition}

\begin{proof}
\emph{Constructing the group and its block.}
Use the actual automorphisms $\alpha_s$ of $\bar N$ and actual
elements $n_{s,t}\in N$ supplied by
Proposition~\ref{prop:structure}(ii).  On the finite set
$\bar N\times S$ define
\begin{equation}\label{eq:extension-group}
 (g,s)(h,t)=\bigl(g\alpha_s(h)n_{s,t},st\bigr).
\end{equation}
Equations~\eqref{eq:actual-actions} and
\eqref{eq:actual-factors} give associativity, and their
normalizations give the identity and inverses.  Hence this is
a finite group $\Gamma_S$ with normal subgroup $\bar N$ and
quotient $S$.  Sending the crossed generator $u_s$ to $(1,s)$
identifies the restriction of the unprojected crossed order with
$\cO\Gamma_S$.

The block $\bar b$ is $S$-stable.  A stable block has a unique
covering block in an extension of $p$-power index.  Its idempotent
is $\bar b$, so our restriction is exactly
\[
 A_S=\cO\Gamma_S\bar b.
\]
\emph{Bounding its defect and coefficient period.}
If $D$ is a defect group of this block, normal block theory gives
a defect group $D\cap\bar N$ of $\bar b$, after conjugacy.
Moreover $D$ maps into $S$, and in this stable $p$-extension it
maps onto $S$.  In particular,
\[
 |D|\leq |S|\,|D\cap\bar N|.
\]
Both quantities on the right are bounded by
Proposition~\ref{prop:structure}.  Notice that the bound is
applied to the group after the central descent in
Section~\ref{sec:comparison}.

Proposition~\ref{prop:comparison}, applied to these same actual
factors, and Lemma~\ref{lem:induced} give a Morita equivalence
between $A_S$ and $\sigma^m A_S$, with $m$ uniformly bounded.
Thus $\mf_{\cO}(A_S)\leq m$.  We now have an actual block with both
required local bounds.  Write $M_1(p,M)$ for the defect-order bound
just obtained.  Proposition~\ref{prop:criteria}(i), applied with
$M_1(p,M)$, bounds its Cartan sum.  Part
(ii) of that proposition gives the claimed finite integral
Morita list.  There are only finitely many abstract groups $S$
of the possible bounded orders.
\end{proof}

\subsection{Full basic corners and arbitrary markings}

Passing to a basic identity component is the integral crossed-product
corner construction of \cite[Proposition~4.15 and Corollary~4.16]{EiseleReduction2021}.
We recall it below to retain the homogeneous units and their labels.

\begin{theorem}[Based finiteness on $p$-subgroups]
\label{thm:based-sylow}
Compress the crossed orders of Proposition~\ref{prop:structure}
by a basic idempotent in the identity block, and identify the
resulting identity order with a representative $R_i$ of its
finite list.  For every possible abstract $p$-group $S$, their
restrictions to $S$ have finitely many based graded isomorphism
classes.  The identity-order identifications may be chosen
arbitrarily over $\cO$.
\end{theorem}

\begin{proof}
\emph{Pure subgroups and a fixed identity order.}
First suppose $S$ is pure.  Let $e$ be a basic idempotent of
$B_0$.  For every homogeneous automorphism, its translate of
$e$ is conjugate to $e$ by a unit of $B_0$: both idempotents
represent a projective module containing one copy of every
indecomposable projective type.  Correcting each homogeneous
unit by such a unit of $B_0$ makes it commute with $e$.
Consequently
\[
 C_S=eA_Se
\]
is a crossed order on $R=eB_0e$, and
\begin{equation}\label{eq:corner-rank}
 \rank_{\cO}C_S=|S|\rank_{\cO}R.
\end{equation}
The idempotent $e$ is full in $A_S$, since it is already full
in $B_0$.

Proposition~\ref{prop:sylow-block} gives finitely many Morita
classes for these total orders.  For fixed $R$ and $S$,
Proposition~\ref{prop:marked} now gives finitely many based
graded types.  The proof of that proposition applies here
because each $C_S$ is Morita equivalent to the actual block
$A_S$.  In particular, its first Hochschild cohomology
vanishes.  Equation~\eqref{eq:corner-rank} also displays
explicitly the rank bound used to pass from Morita classes to
isomorphism classes of total orders.

\emph{Transport from arbitrary subgroups.}
For an arbitrary $p$-subgroup of $Q$, conjugate by a character
in $\Xi$ to make it pure, as explained after
Proposition~\ref{prop:structure}.  Conjugation is implemented
by an integral homogeneous unit and transports the identity
block and its full basic corner.  The transported identity
order is still in the same finite list.  Arbitrary integral
markings are already included in
Proposition~\ref{prop:marked}; inner differences are absorbed
by conjugation in degree one, and outer differences form a
finite set.  Transporting back proves the assertion for the
original subgroup and all its markings.
\end{proof}

The information retained by this theorem is stronger than a Morita
list: an isomorphism fixes the specified copy of $R$ and every group
label.  This is exactly what the kernel argument will need when it
compares central factor systems after correcting the matrix actions.
The integral finiteness came from actual block orders and rigidity,
without a point count over $W(\F_q)$.
